% Generated draft: requires mathematical and source review.
\documentclass[11pt]{article}
\usepackage[margin=25mm]{geometry}
\usepackage{fontspec}
\setmainfont{Latin Modern Roman}
\usepackage{amsmath,amssymb,mathtools,unicode-math}
\setmathfont{Latin Modern Math}
\usepackage{enumitem,needspace,xurl,hyperref}
\setlength{\parindent}{0pt}
\setlength{\parskip}{5pt}
\setlength{\emergencystretch}{3em}
\newcommand{\N}{\mathbb{N}}
\newcommand{\Z}{\mathbb{Z}}
\newcommand{\Q}{\mathbb{Q}}
\newcommand{\R}{\mathbb{R}}
\newcommand{\C}{\mathbb{C}}
\newcommand{\F}{\mathbb{F}}
\newcommand{\A}{\mathbb{A}}
\newcommand{\PP}{\mathbb{P}}
\newcommand{\E}{\mathbb{E}}
\newcommand{\eps}{\varepsilon}
\newcommand{\dd}{\,\mathrm{d}}
\begin{document}
\section*{UnsolvedMath ID: 8800055}
\textbf{Source problem code:} catalogue entry
\subsection{Definitions and mathematical statement}
Let $C/\F_q$ be a smooth projective curve of genus four and let $J=\operatorname{Jac}(C)$, an abelian variety of dimension four. In a usual pairing setting, choose a prime $r$ coprime to $q$ dividing the relevant Jacobian-group order and subgroups used by a nondegenerate bilinear pairing. The embedding degree is
\[
 k=\min\{j\ge1:r\mid q^j-1\},
\]
and the pairing target is contained in $\mu_r\subset\F_{q^k}^\times$. These are standard parameters for describing the tradeoff, not assumptions completely specified by the short source statement.

The task is to determine the exact computational-speed and security tradeoff for genus-four curves used in pairing cryptography. A determinate comparison must fix the curve family, subgroup choices, pairing and arithmetic algorithms, field representation, implementation/cost model, and required security definition. Speed might mean pairing evaluation or other group operations; security must specify attack models, resource costs, and whether estimates or proven bounds are required.

The genus alone does not determine $q$, $r$, $k$, representation, or operation costs. Nor does the source restrict genus-four curves to hyperelliptic curves or prescribe one pairing-friendly construction. It supplies no benchmark platform, parameter range, baseline, optimization objective, or numerical threshold. Hence this is a research-program request for a specified performance/security comparison, not an intrinsic exact function of genus that can be asserted without those choices.
\subsection{Short English statement}
Analyze how genus-four pairing parameters and implementations trade computation cost against the chosen security level. The curve, pairing, algorithm, and security models are necessary parts of the question.

The supplied summary reports no definitive complete tradeoff characterization and calls genus-four pairing-friendly constructions rare, without exact examples, results, or references. Its OPEN-TRIAGE label is retained as a supplied assessment of the 2006 research program; current primary-literature audit remains required.
\subsection{Sources}
\begin{itemize}
\item Local UnsolvedMath registry, permanent ID 8800055, source code catalogue entry, ``Genus-four pairing speed-security tradeoff,'' proposed by Tanja Lange. Dataset curation and metadata: Ulam AI and UnsolvedMath Contributors, CC BY 4.0.
\item IPAM, \emph{Workshop on Number Theory and Cryptography --- Open Problems} (2006), Attachment 6239, PDF page 41, printed slide 39, schedule locator \url{https://www.ipam.ucla.edu/programs/workshops/workshop-i-number-theory-and-cryptography-open-problems/?tab=schedule}; pdf-and-legacy-ppt visual-audit extraction, accessed 2026-07-29.
\item The summary supplies no exact later theorem, construction, or comparison reference.
\item Source provenance reports inspection of 17 PDFs and two legacy PPT decks among 19 surviving attachments, excluding repetitions and solved definitions; this is supplied provenance, not a new inspection claimed here.
\item Current primary-literature status audit and redistribution terms NEEDS\_REVIEW for the publicly accessible author/publisher copy, separately from the dataset licence.
\end{itemize}
\end{document}
